% ============================================================
%  Abstract
% ============================================================
\begin{abstract}
\noindent
本项目研究的是A股日频选股中的\textbf{横截面排序信号}：在每日约4000--5000只可交易股票中，模型需要判断哪些股票相对更值得进入top-20组合。我们不把任务简化为单股次日收益的精确回归，也不声称构建了完整的真实交易系统；本文的核心评估对象是模型在holdout集上的排序能力和选股信号质量，交易成本、滑点、换手约束等执行层因素作为后续交易策略模块处理。基于这一设定，本文提出\textbf{CLIME}（Cross-Modal Injection via Learned Market Encoding），一种面向股票排序的Transformer架构。模型通过两阶段训练先学习pairwise排序先验，再使用方向感知回归损失进行精调；同时，ScaledGatedEncoder将同业动态（peer dynamics）和市场上下文以受控加法方式注入Backbone表征，使市场信息作为增量修正而非直接覆盖个股判断。

在holdout测试集（2025年12月至2026年5月，100个交易日）上，CLIME取得\textbf{+137.15\%累计收益}和\textbf{+123.14\%超额收益}（Sharpe 6.94），比最强基线Ridge（+91.65\%超额）高出约31个百分点，并在全部6个持有月保持正超额。消融实验表明，Stage~1预训练、BCE课程预热、Directional Regression Loss和ScaledGatedEncoder均对性能有明显贡献，其中训练策略组件的影响大于架构组件。除收益对比外，本文还通过Jacobian/SVD、线性探针、特征重要性和极端行情切片分析模型内部行为，验证模型确实利用了横截面排名、基本面和同业上下文等信号，而不是仅依赖单股历史趋势。

最后，本文记录了一次同花顺模拟交易作为端到端压力测试。该模拟最终收益为$-19.40\%$，显著弱于离线结果；但逐日调仓记录显示，主要问题来自日频信号与盘中人工执行之间的协议差异，包括非固定买入时点、先卖后买的资金释放约束、弱势市场下的人工择时压力和后期仓位集中。因此，该结果被解释为\textbf{执行层异常案例}，而不是CLIME离线排序能力的直接否定。它进一步说明，alpha predictor、risk/scoring、portfolio selection和execution必须作为不同模块分别建模。

\end{abstract}
